% generated by roam.py handbook export -- не редактировать, править заметки
% 20261010100652-формула_теилора_с_остаточным_членом_пеано.org (body 527446e6eaae)

\begin{theorem}[{Формула Тейлора с остаточным членом Пеано}]\label[theorem]{note:F7C9E077-9CAB-4149-83B4-67923B49E2F6}
  Пусть $f$ дифференцируема $n$ раз в точке $x_0$. Тогда при $x \to x_0$
  \[
  f(x) = \sum\limits_{k=0}^{n} \frac{f^{(k)}(x_0)}{k!} (x - x_0)^k + o\left((x - x_0)^n\right)
  \]
\end{theorem}

\begin{remark}
  При $n = 1$ это определение дифференцируемости: $f(x) = f(x_0) + f'(x_0)(x - x_0) + o(x - x_0)$~--- именно так формула используется в \hyperref[note:C143295F-C422-4FEE-95AD-3D6B620BC0E8]{теореме о существовании шага для Армихо-Вульфа}.
\end{remark}

\begin{proof}
  \proofpart{Лемма}
  Если $g$ дифференцируема $n$ раз в точке $x_0$ и $g(x_0) = g'(x_0) = \ldots = g^{(n)}(x_0) = 0$, то $g(x) = o((x - x_0)^n)$ при $x \to x_0$.

  Индукция по $n$. База $n = 1$: по определению дифференцируемости $g(x) = g(x_0) + g'(x_0)(x - x_0) + o(x - x_0) = o(x - x_0)$.
  Переход: пусть $n \ge 2$. Функция $g'$ дифференцируема $n - 1$ раз в точке $x_0$, и все её производные до порядка $n - 1$ в $x_0$ равны нулю, поэтому по предположению $g'(t) = o((t - x_0)^{n-1})$: для любого $\eps > 0$ найдётся $\delta > 0$ такое, что $|g'(t)| \le \eps |t - x_0|^{n-1}$ при $|t - x_0| < \delta$.
  Возьмём $x$ с $|x - x_0| < \delta$. По теореме Лагранжа о среднем значении $g(x) = g(x) - g(x_0) = g'(\xi)(x - x_0)$ для некоторого $\xi$ между $x_0$ и $x$. Тогда
  \[
  |g(x)| = |g'(\xi)| \cdot |x - x_0| \le \eps |\xi - x_0|^{n-1} |x - x_0| \le \eps |x - x_0|^n,
  \]
  то есть $g(x) = o((x - x_0)^n)$.

  \proofpart{Теорема}
  Пусть $T(x) = \sum_{k=0}^{n} \frac{f^{(k)}(x_0)}{k!} (x - x_0)^k$~--- многочлен Тейлора и $g = f - T$. Для $m \le n$
  \[
  T^{(m)}(x_0) = \frac{f^{(m)}(x_0)}{m!} \cdot m! = f^{(m)}(x_0),
  \]
  поэтому $g(x_0) = g'(x_0) = \ldots = g^{(n)}(x_0) = 0$, и по лемме $g(x) = o((x - x_0)^n)$.
\end{proof}
